\documentclass[11pt]{article}
\usepackage[margin=1in]{geometry}
\usepackage{amsmath,amssymb,graphicx,booktabs}
\usepackage[authoryear,round]{natbib}
\usepackage{hyperref}

\title{How Volatility Is Delivered: Arias Intensity and Husid Timing\\as Measures of Intraday Market Stress}
\author{Your Name}
\date{\today}

\begin{document}
\maketitle

\begin{abstract}
% Write last. One sentence each: question, method, data, main result, what didn't work.
\end{abstract}

\section{Introduction}
% Question; why timing could matter beyond RV; contributions in order: Husid timing
% diagnostics (headline), bounded liquidity weighting, full-sample event detection.

\section{Related work}
\subsection{Ground-motion intensity and duration}   % Arias 1970; Husid 1969; Trifunac & Brady 1975; Bommer et al. 2009; Afshari & Stewart 2016; Rousseeuw & Leroy 1988 (shortest half)
\subsection{Realized variance and intraday periodicity} % Andersen & Bollerslev 1997; Boudt et al. 2011; Corsi 2009
\subsection{Liquidity measures}                     % Amihud 2002; Corwin & Schultz 2012; Roll 1984
\subsection{Flash events and reversal}              % CFTC-SEC 2010; Kirilenko et al. 2017; Campbell et al. 1993

\section{Data}
% Vendor, cleaning, data-quality report, flagged days.

\section{Method}
\subsection{Returns and deseasonalization}
\subsection{Energy, the Husid curve, and concentration measures}
% Arias intensity on returns = RV, so the level is the benchmark, not the contribution.
% D5-95, D5-75 and W50 as Husid-curve durations; why the start point floats on a trading day
% (constant background energy). Figure: one Husid curve with all three marked.
\subsection{Bounded damping and liquidity weighting}
% Caveat: in seismology damping belongs to the structure, not the ground motion.
\subsection{Event detection}

\section{Pre-registered hypotheses}
% Point to the HYPOTHESES.md commit hash and date.

\section{Results}
\subsection{Simulation: recovering known bursts}
\subsection{H1: concentration and reversal}
\subsection{H2: next-day VIX}
\subsection{H3: differenced signal}

\section{Robustness}

\section{Limitations}

\section{Conclusion}

\bibliographystyle{plainnat}
\bibliography{refs}
\end{document}
